% from ega4-2-fr.tex
\oldpage[IV-2]{42}
soit un $\mathcal O_U$-Module \emph{intègre}~: il suffit en effet de prendre
$U$ tel que $\mathcal F|U$ soit un $\mathcal O_U$-Module réduit (3.2.2).

On dit que le préschéma localement noethérien $X$ est \emph{irredondant} si
$\mathcal O_X$ est irredondant (ce qui implique que $X$ est
\emph{irréductible})~; pour que $X$ soit \emph{intègre}, il faut et il suffit
que le $\mathcal O_X$-Module $\mathcal O_X$ soit intègre ((I, 2.1.8) et
(3.2.1)). Si $\mathcal O_X$ est intègre en un point $x$, c'est-à-dire si
l'anneau $\mathcal O_x$ est intègre, on dit que $X$ est \emph{intègre au point
$x$}. On dit qu'un sous-préschéma fermé $Y$ de $X$ est \emph{primaire dans
$X$} si l'Idéal $\mathcal J$ de $\mathcal O_X$ qui définit $Y$ est primaire
dans $\mathcal O_X$.

\begin{definition}[3.2.5]
\label{IV.2.3.2.5-fr}
Soient $X$ un préschéma localement noethérien, $\mathcal F$ un
$\mathcal O_X$-Module cohérent. On appelle \emph{décomposition irredondante de
$\mathcal F$} une famille $(\mathcal F_\alpha)_{\alpha\in I}$ de
$\mathcal O_X$-Modules quotients de $\mathcal F$ telle que les
$\mathcal F_\alpha$ soient irredondants, que la famille
$(\operatorname{Supp}(\mathcal F_\alpha))$ soit localement finie, et que
l'homomorphisme canonique
$\mathcal F\to\bigoplus_{\alpha\in I}\mathcal F_\alpha$ soit injectif. On dit
qu'une telle décomposition est \emph{réduite} si les ensembles
$\operatorname{Ass}(\mathcal F_\alpha)$ sont deux à deux distincts, et s'il
n'existe aucune partie $J\ne I$ telle que la sous-famille
$(\mathcal F_\alpha)_{\alpha\in J}$ soit une décomposition irredondante de
$\mathcal F$.
\end{definition}

Si $(\mathcal F_\alpha)_{\alpha\in I}$ est une décomposition irredondante
(resp. irredondante réduite) de $\mathcal F$ et si l'on pose
$\mathcal F_\alpha=\mathcal F/\mathcal G_\alpha$, on dit encore que la
famille $(\mathcal G_\alpha)_{\alpha\in I}$ de sous-$\mathcal O_X$-Modules de
$\mathcal F$ est une \emph{décomposition primaire de $0$ dans $\mathcal F$}~;
on observera que l'hypothèse d'injectivité de l'homomorphisme canonique
$\mathcal F\to\bigoplus_{\alpha\in I}\mathcal F_\alpha$ équivaut à la
condition $\bigcap_{\alpha\in I}\mathcal G_\alpha=0$.

Si $(\mathcal F_\alpha)$ est une décomposition irredondante de $\mathcal F$,
dire qu'elle est \emph{réduite} équivaut à dire que les
$\operatorname{Ass}(\mathcal F_\alpha)$ sont deux à deux distincts et contenus
dans $\operatorname{Ass}(\mathcal F)$~; si
$\operatorname{Ass}(\mathcal F_\alpha)=\{x_\alpha\}$ pour tout $\alpha\in I$,
$\alpha\mapsto x_\alpha$ est une \emph{bijection} de $I$ sur
$\operatorname{Ass}(\mathcal F)$~: ces propriétés sont en effet locales et
résultent donc de Bourbaki, \emph{Alg. comm.}, chap.~IV, \S~2,
n\textsuperscript{o}~3, prop.~4.

\begin{proposition}[3.2.6]
\label{IV.2.3.2.6-fr}
Soient $X$ un préschéma localement noethérien, $\mathcal F$ un
$\mathcal O_X$-Module cohérent. Alors il existe une décomposition irredondante
réduite $(\mathcal F^{(x)})_{x\in\operatorname{Ass}(\mathcal F)}$, formée de
$\mathcal O_X$-Modules cohérents tels que pour tout
$x\in\operatorname{Ass}(\mathcal F)$, on ait
$\operatorname{Ass}(\mathcal F^{(x)})=\{x\}$. Pour tout
$x\in\operatorname{Ass}(\mathcal F)$ tel que $\overline{\{x\}}$ ne soit pas
immergé, $\mathcal F^{(x)}$ est déterminé de façon unique comme image de
l'homomorphisme canonique $\mathcal F\to i_*(i^*(\mathcal F))$, où $i$ est le
morphisme canonique $\operatorname{Spec}(k(x))\to X$.
\end{proposition}

Pour tout $x\in\operatorname{Ass}(\mathcal F)$, soit $U$ un voisinage ouvert
affine de $x$, d'anneau $A$, et soit $\mathcal F|U=\widetilde M$, où $M$ est
un $A$-module de type fini. On sait (Bourbaki, \emph{Alg. comm.}, chap.~IV,
\S~1, n\textsuperscript{o}~1, prop.~4) qu'il existe un sous-module $N$ de $M$
tel que, si l'on pose $P=M/N$, on ait
$\operatorname{Ass}(P)=\{x\}$ et
$\operatorname{Ass}(N)=\operatorname{Ass}(M)-\{x\}$. Soit
$\mathcal G=\widetilde P$, qui est un $\mathcal O_U$-Module quasi-cohérent, et
soit $j$ l'injection canonique $U\to X$~; soit
$u:j^*(\mathcal F)\to\mathcal G$ l'homomorphisme surjectif correspondant à
l'homomorphisme $M\to P$~; on en déduit un homomorphisme
$j_*(u):j_*(j^*(\mathcal F))\to j_*(\mathcal G)$, d'où par composition un
homomorphisme
\[
v:\mathcal F\xrightarrow{\rho_{\mathcal F}}j_*(j^*(\mathcal F))
\xrightarrow{j_*(u)}j_*(\mathcal G)
\]
dont $u$ est la restriction à $U$~; nous désignerons par
$\mathcal F^{(x)}$ l'image de $\mathcal F$ par cet homomorphisme, qui est un
$\mathcal O_X$-Module cohérent (I, 6.1.1). On a
$\operatorname{Ass}(j_*(\mathcal G))=\{x\}$ en vertu de (3.1.13), et
\emph{a fortiori} (3.1.7)
$\operatorname{Ass}(\mathcal F^{(x)})=\{x\}$, puisque
$\mathcal F^{(x)}\ne0$. En outre, si

